\documentclass[11pt]{article}
\usepackage[T1]{fontenc}
\usepackage{amsmath,amssymb,amsthm}
\usepackage[margin=1in]{geometry}
\usepackage[hidelinks]{hyperref}
\newtheorem{theorem}{Theorem}[section]
\newtheorem{proposition}[theorem]{Proposition}
\newtheorem{corollary}[theorem]{Corollary}
\theoremstyle{definition}
\newtheorem{definition}[theorem]{Definition}
\newtheorem{example}[theorem]{Example}
\newtheorem{question}[theorem]{Question}
\newcommand{\Sh}{\operatorname{Sh}}
\title{Cube Bases and Complements of Cube-Chordal Ample Classes}
\author{}
\date{Research notes, 19 October 2026}
\begin{document}
\maketitle

\section{The complementary existence problem}
Let $E$ be a finite labelled coordinate set and $W=\{0,1\}^E$.
All complements below are taken in $W$, retaining original constants.
Write $D(u,v)=\{i:u_i\ne v_i\}$ and
$Q(v,T)=\{w:w|_{E\setminus T}=v|_{E\setminus T}\}$.
For $q_i\in(0,1)$ set
\[
 q_T=\prod_{i\in T}q_i,\qquad t_T=\prod_{i\in T}(1-q_i^2),
 \qquad (G_H)_{uv}=q_{D(u,v)}\quad(u,v\in H).
\]
The actual rooted-cube column is
\[
 p_{v,T}(w)=\begin{cases}(-1)^{|D(v,w)|}q_{D(v,w)},&w\in Q(v,T),\\
                    0,&w\notin Q(v,T),\end{cases}
 \qquad Q(v,T)\subseteq H.
\]
An ordinary map chooses one such support $r(v)$ at each $v\in H$ and
satisfies
\begin{equation}\label{eq:ordinary}
 \sum_{v\in Q}\mathbf1\{\operatorname{dir}(Q)\subseteq r(v)\}=1
 \quad\text{for every positioned cube }Q\subseteq H.
\end{equation}
Cycles are permitted. A support $T$ is shattered if the projection onto
$T$ is all of $\{0,1\}^T$; $\Sh(H)$ denotes the family of such supports.
The class is ample if its shattered supports are
exactly those supports of contained positioned cubes. The class is in
$\mathrm{Ext}$ when both it and its original ambient complement have
ordinary maps.
We use the classical ample-class facts $|H|=|\Sh(H)|$, closure under
conditioning, projection, reduction and ambient complementation, and
isometry of nonempty classes. Counting and coordinate-closure proofs are
recalled in the \href{../sections/ample_counts.tex}{first paper},
\texttt{prop:sandwich}; the isometry formulation is also recorded in
the \href{../../minor-atlas/research/ample_structure.md}{ample-structure module}.

For a nonempty class, form the intersection graph of all inclusion-maximal
positioned cubes, joining cubes even when their intersection is a single
word. A class is in $\mathrm{CAmp}$ if it is ample and this graph is
chordal, meaning it has no induced cycle of length at least four.
A nonempty class $D$ is a daisy rooted at $z\in D$ if
$Q(z,D(z,v))\subseteq D$ for every $v\in D$.
For $v\in D$ put $N_D(v)=\{i:v\text{ with bit }i\text{ flipped lies in }D\}$.
A corner is a word $v$ with $Q(v,N_D(v))\subseteq D$.
A corner peeling is an order in which every word is a corner of the
class remaining at its deletion; the final singleton may be retained.
The empty class has the empty order.
The main result below proves $\mathrm{CAmp}\subseteq\mathrm{Ext}$.
In fact the full original ambient complement has a corner peeling.
The primal corner peeling is already known and is reused. The new
construction enlarges a tree of daisy pieces to one daisy by clamped
cube attachments, then transfers complementary cube columns backwards.

The atlas incremental construction supplies a cube $C=Q(z,S)$ attached to
an old ample class $B$. Put $Y=E\setminus S$, $F=B\cap C$ and
$A=C\setminus B$. Its clamp condition is
\begin{equation}\label{eq:clamp}
 c(b)=(b|_S,z|_Y)\in B\qquad(b\in B).
\end{equation}
The old maximal cubes meeting $C$ form a clique. Consequently $F$ is a
nonempty proper daisy cube rooted at a common word $z$:
$Q(z,D(z,f))\subseteq F$ for every $f\in F$.
The original generator and its proof are in the
\href{../../minor-atlas/sections/cube_chordal_ample.tex}{atlas module},
\texttt{prop:cube-chordal-generator}. They are inputs, not new assertions here.

The next section gives a relative integral congruence. It preserves the
old vertex module, the attachment cube module and its subcubes containing
$z$, and uses actual cube columns at every new root. It does not fix
every singleton subspace: a nontrivial cube column necessarily has several
nonzero vertex coefficients. The allowed support and root of each column
remain part of the data throughout.

\section{Canonical elimination across a daisy interface}
\label{sec:relative}
For this section $B$ need not be ample. Assume \eqref{eq:clamp}, assume
$F$ is a nonempty daisy cube rooted at $z$, and put $H=B\cup C$.
For each $a\in A$ set $T_a=D(z,a)\subseteq S$. Form a square matrix $R$
on the vertex module of $H$: at old roots $b\in B$ its column is the
singleton vector $\delta_b$, and at $a\in A$ its column is $p_{a,T_a}$.

\begin{theorem}[Relative integral cube elimination]\label{thm:relative}
With the same row and root ordering,
\begin{equation}\label{eq:relative}
 R^tG_HR=G_B\oplus\operatorname{diag}_{a\in A}t_{T_a},
 \qquad \det R=1
\end{equation}
over $\mathbb Z[q_i:i\in E]$. Every column is an actual rooted-cube
column. The map $R$ fixes the vertex module on $B$ pointwise and preserves
the module on $C$ and on every $Q(z,T)$ with $T\subseteq S$.

At interior parameters, Schur complementation on the old vertices gives
\begin{equation}\label{eq:schur}
 G_{AA}-G_{AB}G_B^{-1}G_{BA}
   =G_{AA}-G_{AF}G_F^{-1}G_{FA}.
\end{equation}
Thus the conditional form of the new vertices is determined by the actual
cube $C$ and its daisy interface $F$.
\end{theorem}
\begin{proof}
Reflect bits so that $z$ is zero in $S$, solely for describing subset
inclusion. The words of $F$ form a down-set in $2^S$, and the words of
$A$ its complementary up-set in $C$. Each $p_{a,T_a}$ is supported on
the full lower cube from $z$ to $a$, which is contained in $C$.

On a full cube the columns $p_{a,D(z,a)}$, over all roots, are mutually
$G_C$-orthogonal with squared norms $t_{D(z,a)}$, by tensoring the
two-by-two identity. Moreover $(G_Cp_{a,T_a})(f)$ vanishes unless $f$
agrees with $a$ on $T_a$. No $f\in F$ can have that agreement when
$a\notin F$, by downward closure. Hence $G_{F,C}p_{a,T_a}=0$.

For $b\in B$ and $v\in C$, coordinate separation gives
\[
 (G_H)_{bv}=q_{D(b,c(b))}(G_C)_{c(b),v}.
\]
The clamp belongs to $F$, so $G_{B,H}p_{a,T_a}=0$. The old columns
have Gram matrix $G_B$, and the new columns are mutually orthogonal,
proving the congruence in \eqref{eq:relative}.

Order the roots of $A$ by increasing $|T_a|$. The $A,A$ block of $R$
has nonzero entries only at lower subsets of its column root, and has
diagonal one. It is triangular with determinant one. The $A,B$ block
is zero and the $B,B$ block is identity, so $\det R=1$. This proof is
polynomial; no inverse or interior specialization is needed for that claim.
The same triangular argument on $C$, or on a lower cube $Q(z,T)$,
shows that its vertex module is preserved. The inverse restrictions are
polynomial because each corresponding determinant is one.

For the Schur assertion, the same clamp factorization says, row by row,
\[
 G_{BF}G_F^{-1}G_{FA}=G_{BA}.
\]
Let $J$ inject the $F$-coordinates into the $B$-coordinates. Then
$G_BJ=G_{BF}$, whence $G_B^{-1}G_{BA}=JG_F^{-1}G_{FA}$.
Substitution proves \eqref{eq:schur}. The relevant principal matrices
are positive definite at interior parameters.
\end{proof}

If $B$ is ample, the clamp condition makes $H$ ample by the atlas proof,
even before its chordality hypotheses are imposed. Consequently any
ordinary map on $B$ extends to $H$ by the fixed new supports $T_a$.
Indeed substitute its selected cube matrix for the singleton old columns
in \eqref{eq:relative}; all resulting columns are actual and orthogonal,
so the selected-basis criterion in the
\href{../main.pdf}{completed first paper} applies. The combinatorial
extension was already available by deleting $A$ in decreasing subset order
and reversing the corner order. Formula \eqref{eq:relative} makes the
relative form and the preserved attachment modules explicit.

\begin{corollary}[Inverse and trace under attachment]\label{cor:relative-trace}
Under the assumptions of Theorem~\ref{thm:relative}, embed $G_B^{-1}$
in the vertex module of $H$ by zero outside $B$. Then
\begin{equation}\label{eq:inverse-update}
 G_H^{-1}=\widetilde{G_B^{-1}}+
          \sum_{a\in A}\frac{p_{a,T_a}p_{a,T_a}^{t}}{t_{T_a}}.
\end{equation}
In particular the inverse entries outside $C\times C$ are unchanged, and
\begin{equation}\label{eq:trace-update}
 \operatorname{tr}G_H^{-1}-\operatorname{tr}G_B^{-1}
   =\sum_{a\in A}\prod_{i\in T_a}\frac{1+q_i^2}{1-q_i^2}.
\end{equation}
Neither assertion requires ampleness of $B$.
\end{corollary}
\begin{proof}
Invert \eqref{eq:relative} to obtain
$G_H^{-1}=R\operatorname{diag}(G_B^{-1},(t_{T_a}^{-1})_{a\in A})R^t$.
The old columns of $R$ are singleton vectors, which gives
\eqref{eq:inverse-update}. Every added column is supported in $C$.
Its squared Euclidean norm is
$\sum_{w\in Q(a,T_a)}q_{D(a,w)}^2=\prod_{i\in T_a}(1+q_i^2)$.
Taking traces proves \eqref{eq:trace-update}.
\end{proof}

The trace increment depends only on the new supports $T_a$.
The inverse update also retains their positions through the vectors
$p_{a,T_a}$. This is the distinction needed when compatibility depends
on where attachment cubes meet.

\section{The dual product alphabet}
\label{sec:dual}
We now assume $B$ is ample. Work in the physical space
$\mathcal V=(\mathbb R^2)^{\otimes E}$. In coordinate $i$ put
\[
 g_0=(1,0)^t,\quad g_1=(q_i,\sqrt{1-q_i^2})^t,
 \quad e_{00}=g_0,\quad e_{10}=g_1,
 \quad e_{01}=(\sqrt{1-q_i^2},-q_i)^t,\quad e_{11}=(0,1)^t.
\]
For a word $v$ and support $T$, set
$g_v=\bigotimes_i g_{v_i}$ and
$e_{v,T}=\bigotimes_i e_{v_i,\mathbf1\{i\in T\}}$.
Let $V_D=\operatorname{span}\{g_v:v\in D\}$ for a binary class $D$.
The full family $(g_v)_{v\in W}$ is a basis and has Gram matrix $G_W$.
The first-paper dictionary gives
\begin{equation}\label{eq:cube-dictionary}
 \sum_w p_{v,T}(w)g_w=\sqrt{t_T}\,e_{v,T}.
\end{equation}

\begin{proposition}[Complementary product vectors]\label{prop:dual}
For every binary class $D$, not necessarily ample,
\begin{equation}\label{eq:admissible-dual}
 e_{v,T}\in V_D^\perp
 \quad\Longleftrightarrow\quad Q(v,E\setminus T)\subseteq D^c.
\end{equation}
There is an orthogonal involution $U$ on $\mathcal V$ with
\begin{equation}\label{eq:dual-involution}
 UV_{D^c}=V_D^\perp,\qquad
 Ue_{v,T}=(-1)^{|v|}e_{v,E\setminus T}.
\end{equation}
Here $|v|$ counts one bits on the full original frame.
\end{proposition}
\begin{proof}
In a coordinate of $T$, the inner product $\langle g_{h_i},e_{v_i,1}\rangle$
is zero exactly when $h_i\ne v_i$. In other coordinates all local inner
products are nonzero. Thus $\langle g_h,e_{v,T}\rangle\ne0$ precisely
when $h|_T=v|_T$. There is no such $h\in D$ exactly when the cylinder
$Q(v,E\setminus T)$ lies in $D^c$, proving \eqref{eq:admissible-dual}.

Take $U=\bigotimes_i U_i$, where
\[
 U_i=\begin{pmatrix}\sqrt{1-q_i^2}&-q_i\\-q_i&-\sqrt{1-q_i^2}\end{pmatrix}.
\]
Each $U_i$ is an orthogonal involution and sends
$e_{v_i,s}$ to $(-1)^{v_i}e_{v_i,1-s}$. Moreover
$\langle g_u,Ug_v\rangle$ is zero unless $u=v$, when it equals
$(-1)^{|v|}\sqrt{t_E}$. Hence $UV_{D^c}\subseteq V_D^\perp$;
their dimensions are both $|D^c|$, proving equality.
\end{proof}

Thus complementary ordinary existence asks for an orthonormal basis of
$V_H^\perp$ from the prescribed product alphabet, with each root of
$H^c$ used once. Membership of these products is exactly actual cube
membership by \eqref{eq:admissible-dual}, not generic tensor positivity.
The involution uses square roots; we do not assert it is an integral
polynomial congruence. The integral statement is \eqref{eq:relative}.

\begin{proposition}[Clamps and commuting projections]\label{prop:commuting-clamp}
Let $B\subseteq W$ be arbitrary and let $C=Q(z,S)$. Write $\Pi_D$
for the orthogonal projection onto $V_D$. Then
\[
 \Pi_B\Pi_C=\Pi_C\Pi_B
 \quad\Longleftrightarrow\quad c(B)\subseteq B,
 \qquad c(b)=(b|_S,z|_{E\setminus S}).
\]
When these conditions hold, with $F=B\cap C$,
\begin{equation}\label{eq:commuting-union}
 \Pi_B\Pi_C=\Pi_F,\qquad
 \Pi_{B\cup C}=\Pi_B+\Pi_C-\Pi_F,\qquad
 I-\Pi_{B\cup C}=(I-\Pi_B)(I-\Pi_C).
\end{equation}
No ampleness or daisy assumption is needed.
\end{proposition}
\begin{proof}
The tensor formula for projection onto $V_C$ gives
$\Pi_Cg_b=q_{D(b,c(b))}g_{c(b)}$. Every coefficient is nonzero in
the open parameter box. Since $(g_w)_{w\in W}$ is a basis,
$\Pi_CV_B\subseteq V_B$ holds exactly when $c(B)\subseteq B$.
For two orthogonal projections, invariance of one range under the other
is equivalent to commutation: self-adjointness also makes the orthogonal
complement invariant, giving a block diagonal operator.
The product of commuting orthogonal projections projects onto the
intersection of their ranges. The full vertex basis implies
$V_B\cap V_C=V_{B\cap C}$. The sum and complementary product in
\eqref{eq:commuting-union} follow from the simultaneous range/kernel
decomposition of the two commuting projections.
\end{proof}

\section{A bijection with a prescribed old-complement fiber}
\label{sec:excision}
Put $K_0=B^c$ and $K=H^c=K_0\setminus A$.
Both $K_0$ and $K$ are ample, by complement closure and
the ample union assertion above; this does not assert ordinary existence.
The clamp has a useful
consequence: if an $S$-trace is absent from $F$, its entire $Y$-fiber
is absent from $B$. Thus $A\times\{0,1\}^Y$, interpreted by its
$S$-traces, lies in $K_0$. Since $A$ is an up-set relative to $z$,
\begin{equation}\label{eq:pinned-cubes}
 Q(a,E\setminus T_a)\subseteq K_0\qquad(a\in A).
\end{equation}
These are actual full original cubes, including all $Y$-directions.

\begin{theorem}[Exact complementary excision fiber]\label{thm:excision}
Ordinary maps on $K$ are in bijection with ordinary maps $r_0$ on $K_0$
satisfying the simultaneous prescribed supports
\begin{equation}\label{eq:pins}
 r_0(a)=E\setminus T_a\qquad\text{for every }a\in A.
\end{equation}
The bijection retains every rooted-cube column at roots in $K$ and adds
exactly the columns \eqref{eq:pinned-cubes} at the removed roots.
In particular ordinary existence on $K$ is equivalent to nonemptiness of
this one prescribed-map fiber on $K_0$.
\end{theorem}
\begin{proof}
By \eqref{eq:relative} and \eqref{eq:cube-dictionary}, the unit vectors
$e_{a,T_a}$ are mutually orthogonal and orthogonal to $V_B$. Their span
is the orthogonal complement of $V_B$ inside $V_H$, by the dimension
count $|H|-|B|=|A|$. Therefore
\begin{equation}\label{eq:orthogonal-excision}
 V_B^\perp=V_H^\perp\oplus
                 \operatorname{span}\{e_{a,T_a}:a\in A\}.
\end{equation}

Let $r$ be an ordinary map on $K$. Its unit product vectors are an
orthonormal basis of $V_K$. Apply $U$ and disregard rootwise signs to
obtain an orthonormal basis of $V_H^\perp$. Add the vectors $e_{a,T_a}$
in \eqref{eq:orthogonal-excision}. Applying $U$ back gives the unit
product vectors selected at roots $K$ by $r$ and at roots $a\in A$ by
\eqref{eq:pins}. All the latter cubes are actual by
\eqref{eq:pinned-cubes}. The ordinary-map/orthogonal-basis equivalence
on the ample class $K_0$ proves that this is an ordinary map there.

Conversely take an ordinary $r_0$ on $K_0$ satisfying \eqref{eq:pins}.
After applying $U$, its remaining product vectors span $V_H^\perp$ by
\eqref{eq:orthogonal-excision}. Proposition~\ref{prop:dual} implies that
each corresponding unchanged rooted cube is wholly in $K$. The remaining
roots are exactly $K$, and their vectors are orthonormal. The same
selected-basis equivalence on the ample class $K$ proves ordinariness.
Both directions fix each column at each retained root, so are inverse.
\end{proof}

\begin{example}[A noncube interface with several removed roots]
Take $S=\{1,2,3\}$, $Y=\{4\}$, and
$F=\{000,100,010,001\}$ in the $S$-coordinates. Let
$B=F\times\{0,1\}$ and $C=\{0,1\}^S\times\{0\}$, with root $z=0000$.
The four roots in $A$ have $S$-supports $12,13,23,123$. The exact old
complement supports prescribed by \eqref{eq:pins} are, respectively,
$\{3,4\},\{2,4\},\{1,4\},\{4\}$. All four constraints must hold in
one whole ordinary map. Separate feasibility of these actual columns
does not establish simultaneous feasibility.
\end{example}

\begin{question}\label{qu:supply}
For the simplicial-cube attachments in the atlas generator, can one
supply an ordinary map on $B^c$ satisfying all of \eqref{eq:pins}?
\end{question}
Inductively supplying this fiber would prove $\mathrm{CAmp}\subseteq\mathrm{Ext}$.
The relative congruence alone proves neither its nonemptiness nor preservation
of an arbitrarily preselected old-complement map. The next section supplies
an opposite by enlarging the whole class and using excision in the reverse
direction. It does not require a preselected old-complement map to survive.

\section{Completing a tree of daisies by cube attachments}
\label{sec:tree-completion}
The useful induction runs towards a larger class whose opposite is easy.
Its pieces are allowed to change. We preserve a tree of rooted daisy
pieces, rather than assume the enlarged class still has a chordal
maximal-cube graph.

Here is the precise tree input. Let $T$ be a finite tree with nonempty
pieces $D_t\subseteq W$ and roots $z_t$. Require:
\begin{enumerate}
\item each $D_t$ is a daisy rooted at $z_t$;
\item for every nonempty connected subtree $U$, the union
      $H_U=\bigcup_{t\in U}D_t$ is ample, and every positioned cube in
      $H_U$ is contained in a piece $D_t$ with $t\in U$;
\item for each word $v\in H_T$, the nodes $\{t:v\in D_t\}$ form
      a connected subtree;
\item each adjacent overlap $F_{tu}=D_t\cap D_u$ is nonempty and
      is a daisy rooted at both $z_t$ and $z_u$.
\end{enumerate}
Every nonempty cube-chordal ample class has such a tree. This is the
forward, subtree and overlap part of the previously established
\href{../../maxmin-forbidden/sections/chordal_downset_tree.tex}
{positioned tree certificate},
\texttt{thm:maxmin-chordal-downset-tree}; its complete proof is the
input here, not an assertion of compatible opposite maps. Its pieces
are unions of the maximal cubes in clique bags. A common word roots
each bag; connected label occurrences give item~3 and the exact edge
overlaps. Connected subunions retain ampleness and their exact cube
labels by the simplicial-cube removal argument. The shared cube labels
contain both adjacent roots, proving item~4. No maximum-parent or
ordinary complementary existence claim from that theorem is used.

\begin{proposition}[Aligning the root of a leaf piece]
\label{prop:leaf-root-alignment}
Suppose $B$ and $H=B\cup D$ are ample, $D$ is a nonempty daisy rooted
at $z$, and $F=B\cap D$ is a nonempty daisy rooted at both $z$ and $w$.
Suppose every positioned cube of $H$ is contained in $B$ or $D$.
Set $J=D(z,w)$ and
\[
 D^+=\{x\in W:x|_{E\setminus J}\in\pi_{E\setminus J}D\}.
\]
Then $D^+$ is a daisy rooted at $w$, and
\begin{equation}\label{eq:private-leaf-cylinder}
                  D^+\cap B=F.
\end{equation}
There is a finite sequence from $H$ to $B\cup D^+$ consisting of
clamp-preserving cube attachments with nonempty daisy interfaces.
Every attachment root is $z$. Each intermediate class is ample and
every newly created cube is contained in the cube just attached.
\end{proposition}
\begin{proof}
First $F$ is invariant under changing arbitrary $J$-coordinates. To
see this, flip coordinates so that $z=0$ and $w$ is the indicator of
$J$. For $f\in F$, downward closure towards $z$ puts $f\setminus J$
in $F$. The interval from $w$ to $f\setminus J$, contained in $F$
by the other daisy orientation, contains every $J$-assignment with
outside trace $f|_{E\setminus J}$. This proves the invariance.
Also $D^+=\{0,1\}^J\times\pi_{E\setminus J}D$ is a daisy at every
root with outside trace $z|_{E\setminus J}$, including $w$.

We prove \eqref{eq:private-leaf-cylinder}. Take $b\in B\cap D^+$,
and let $P$ be the full $J$-fiber through $b$. Some $d\in D$ lies in
$P$. If $P$ already meets $F$, its whole fiber lies in $F$ by the
invariance just proved. Otherwise $d\in D\setminus B$ and
$b\in B\setminus D$. The nonempty conditioning $H\cap P$ is ample,
hence connected. A vertex path from $d$ to $b$ must meet $F$:
every edge of $H$ is contained in $B$ or $D$, so no edge can join
$D\setminus B$ directly to $B\setminus D$. This contradiction proves
$b\in F$. The reverse containment is immediate.

For a rooted cube $C=Q(z,S)\subseteq D$, the clamp of $B$ into $C$
lands in $F\cap C$. We give the needed compatibility argument.
Cube carriage and ampleness imply
\[
 \Sh(H)=\Sh(B)\cup\Sh(D),\qquad
 \Sh(B)\cap\Sh(D)=\Sh(F).
\]
For the second equality, its right side is contained in its left side,
and the ample cardinalities and $|H|=|B|+|D|-|F|$ make their sizes
equal. Since $2^S\subseteq\Sh(D)$, projection closure and the same
cardinality identity give
\[
 |\pi_SB|=|\Sh(B)\cap2^S|=|\Sh(F)\cap2^S|=|\pi_SF|.
\]
The inclusion $\pi_SF\subseteq\pi_SB$ is therefore equality.
Clamping a word of $F$ into $Q(z,S)$ stays in $F$, by downward closure
towards $z$. Matching the $S$-trace proves the clamp assertion for $B$.

List the finitely many maximal rooted cubes $C_i=Q(z,S_i)$ of $D$,
and attach $C_i^+=Q(z,S_i\cup J)$ successively, skipping empty changes.
Their union is $D^+$. For $b\in B$, its clamp into $C_i^+$ differs
only in $J$ from its clamp into $C_i$, so lies in $F$.
Clamping a word of $D$ into $C_i^+$ moves coordinates towards $z$,
and therefore stays in $D$. The same is true of every previously
attached $C_j^+$, which also contains $z$. Thus each current class
is clamp-invariant for the next cube.
Write $D_i=D\cup\bigcup_{j<i}C_j^+$, a daisy rooted at $z$.
Equation~\eqref{eq:private-leaf-cylinder} gives
\[
 (B\cup D_i)\cap C_i^+=D_i\cap C_i^+,
\]
a nonempty daisy rooted at $z$. The ample-union and no-mixed-cube
argument of the atlas clamp theorem now applies at every step.
This proves all assertions, including the case $J=\varnothing$.
\end{proof}

\begin{theorem}[Complementary corner peelings for cube-chordal classes]
\label{thm:camp-ext}
Every nonempty cube-chordal ample class $H\subseteq W$ and its full
original ambient complement have corner peelings. Empty complements
use the empty order. In particular
\[
                         \mathrm{CAmp}\subseteq\mathrm{Ext}.
\]
More generally the conclusion holds for any union having the four
tree-of-daisies properties above. An ordinary complementary map can be
constructed using only the original bag roots: for every $a\in H^c$,
\begin{equation}\label{eq:tree-opposite}
                 r_{H^c}(a)=E\setminus D(a,z(a)),
\end{equation}
where $z(a)$ is selected by the finite completion below. It is not a
first-eligible-root rule. Every original constant coordinate is retained.
\end{theorem}
\begin{proof}
We first complete the union to a single daisy while preserving the tree
properties. We then peel the complement by undoing the completion.

\emph{Leaf alignment and the induction invariant.}
Choose a leaf $t$ with parent $u$, set $D=D_t$ and
$B=\bigcup_{v\ne t}D_v$. Word connectivity gives
$D\cap B=D_t\cap D_u=F$. The connected-subtree assumption makes
$B,H$ ample, and every cube of $H$ belongs to $B$ or $D$.
Proposition~\ref{prop:leaf-root-alignment} applies with
$z=z_t$, $w=z_u$. Expand $D$ to $D^+$ by its stated cube attachments,
and replace $D_u$ by $D_u\cup D^+$, deleting leaf $t$.
The new parent piece is a daisy rooted at $z_u$, since both terms
are daisies with that root. Equation~\eqref{eq:private-leaf-cylinder}
shows, for every other retained piece $D_v$,
\[
 (D_u\cup D^+)\cap D_v=D_u\cap D_v.
\]
All remaining adjacent overlaps are consequently unchanged, with
their two original roots.

Word connectivity also persists. An old word exclusive to $D_t$
now occurs only at $u$. Any other old word occurring at $t$ already
occurs at $u$, by connectivity; delete $t$ from its occurrence subtree.
Every genuinely new word occurs only in the new parent piece, because
$D^+\cap B=F$. All other occurrence subtrees are unchanged.

It is necessary to check every connected subunion, not only the global
one. A new subtree avoiding $u$ is unchanged. For a new connected
subtree $U$ containing $u$, put $B_U=\bigcup_{v\in U}D_v$, using the
old pieces. The old subtree $U\cup\{t\}$ is connected, so
$B_U\cup D$ is ample and its cubes are carried by its pieces.
In the proof of Proposition~\ref{prop:leaf-root-alignment}, every
clamp of a word of $B$ into an expanded cube lies in $F$.
Thus clamps of $B_U$ land in $F\subseteq D_u\subseteq B_U$.
The same sequence of cube attachments is therefore valid on
$B_U\cup D$. It preserves ampleness, and each new cube is carried
by an expanded leaf cube, hence by the merged parent piece.
Old cubes remain carried as well. This proves item~2 for the new tree.
Items~1, 3 and 4 were verified above, so the entire induction invariant
is preserved. Each merge strictly decreases the number of pieces.

We obtain a finite increasing sequence
\begin{equation}\label{eq:tree-completion-chain}
 H=H_0\subset H_1\subset\cdots\subset H_m=L,
\end{equation}
where $L$ is one daisy rooted at a remaining original bag root $z_*$.
Each strict step attaches a positioned cube with a nonempty daisy
interface and a clamp preserving the old class. Denote its root by
$z_j$ and its new words by $A_j=H_{j+1}\setminus H_j$.
Every $z_j$ is an original bag root.

\emph{The complementary order and its whole outgoing cubes.}
For one step write $C=Q(z,S)$, $F=H_j\cap C$ and $A=C\setminus H_j$.
Inside $C$, $A$ is an up-set relative to $z$.
In $H_j^c$, remove $A$ in increasing $|D(z,a)|$, with arbitrary ties.
At the deletion of $a$, its remaining incident directions are exactly
\[
                         E\setminus D(z,a).
\]
Indeed a direction in $D(z,a)\subseteq S$ moves to a smaller word of
$C$, which is either in $F$ or was removed earlier. A direction in
$S\setminus D(z,a)$ moves to a larger word of $A$, still present.
Every direction outside $S$ leads outside $C$ and is also present:
an old word there with the same $S$-trace would clamp to the missing
word $a$, contradicting clamp invariance.
The whole cube on those directions remains present. Each of its
$S$-traces has difference support containing $D(z,a)$, so belongs to $A$;
all parallel outside assignments
are absent from $H_j$ by the same clamp argument. Points in $C$
are still unremoved words of $A$, and points outside $C$ belong to
$H_{j+1}^c$. Thus $a$ is a corner, not merely the owner of its edges.
Removing this block leaves exactly $H_{j+1}^c$.

Process $A_0,A_1,\ldots,A_{m-1}$ in that order. The remaining class
is $L^c$, an up-set relative to $z_*$; it peels by increasing distance
from $z_*$, with outgoing support $E\setminus D(a,z_*)$.
This proves the complementary corner assertion and
\eqref{eq:tree-opposite}, taking $z(a)=z_j$ on $A_j$ and $z_*$
on $L^c$. A last surviving word may be retained rather than deleted.
Equivalently, Theorem~\ref{thm:excision} transfers the canonical ordinary
map on $L^c$ backwards across every step, retaining each completed
column and supplying all new pins simultaneously.

\emph{The primal order.}
The primal peeling is the established tree peeling, whose short argument
also applies to the four stated properties. In the original tree remove
a leaf piece $D_t$ while retaining the other union $B$.
Its exclusive words $D_t\setminus B$ form an up-set in the daisy
$D_t$ relative to $z_t$, since the overlap is a daisy at $z_t$.
Delete them in decreasing $|D(a,z_t)|$.
Every edge at an exclusive word is carried by $D_t$; no cube in a
different piece contains that word. Rootward neighbors and their full
cube remain. Any opposite neighbor is another exclusive word of larger
distance, already removed. Its incident cube is therefore exactly
$Q(a,D(a,z_t))$. Delete the leaf and repeat, finally peeling the last
daisy in decreasing distance from its root. This gives the primal
ordinary map and completes both-shore existence.
\end{proof}

If the empty class is admitted into $\mathrm{CAmp}$ by convention,
the same conclusion holds: its primal map is empty and its complement
is the full cube with its standard corner order.

\begin{corollary}[An integral actual-cube certificate of completion]
\label{cor:tree-integral-completion}
For the chain \eqref{eq:tree-completion-chain}, put $z(a)=z_j$
when $a\in A_j$. On the vertex module of $L$, let $R$ have old columns
$\delta_h$ for $h\in H$ and column $p_{a,D(a,z(a))}$ for
$a\in L\setminus H$. Then
\[
 R^tG_LR=G_H\oplus
       \operatorname{diag}_{a\in L\setminus H}t_{D(a,z(a))},
 \qquad\det R=1.
\]
Every column is an actual rooted-cube column, and $R$ fixes the old
vertex module pointwise. Replacing the old singleton bank by any
supplied ordinary map on $H$ gives an ordinary cube basis on $L$.
Together with the canonical complementary bank outside $L$, this also
constructs a full ambient native basis extending that supplied primal
map. Its polynomial actual-cube matrix is unimodular.
\end{corollary}
\begin{proof}
Apply Theorem~\ref{thm:relative} successively, extending each earlier
matrix by identity at later roots. Every later matrix fixes the entire
earlier vertex module, so all previously constructed actual columns
remain literally unchanged. This gives the displayed matrix, congruence
and determinant. Normalizing the added columns gives an orthonormal
basis of $V_L\cap V_H^\perp$. Replace the old bank by the supplied
orthonormal primal bank, and complete by the dual of the canonical
ordinary bank on $L^c$. The selected-basis criterion on the full cube
gives its polynomial unimodularity.
\end{proof}

Consequently Question~\ref{qu:supply} has an affirmative answer for
every atlas cube-chordal attachment. Apply Theorem~\ref{thm:camp-ext}
to the new class $B\cup C$ and then Theorem~\ref{thm:excision}:
the constructed new-complement map supplies all canonical pins on $B^c$
simultaneously. Any supplied old primal map can be retained by
Theorem~\ref{thm:relative}. No independently chosen local bases are
assumed to agree.

The independent exact construction in \texttt{verify\_tree\_root\_alignment.py}
checks four fixed existing inputs: a cube with original constants, the
non-Horn atlas class $V$, the chordal sun, and the alternating-root path.
It verifies all intermediate clamps and daisy interfaces, every connected
subunion in these fixed trees, cube carriage, both final ordinary maps,
and the complementary corner orders. There are 52,281 positioned-equation
checks, 82 connected-subunion checks and 3,942 cube-carriage checks.
The relative matrices are triangular with diagonal entries one, and all 65,896
full ambient native-vector pairs are checked exactly.
These tests support implementation correctness; the arbitrary-dimensional
proof is the invariant and induction above. The tree certificate and
primal peeling are old inputs. The additional step is leaf root alignment
with private-cylinder overlap, followed by complementary excision.

\section{Closure intervals give complementary cube bases}
\label{sec:closure-completion}
Suppose $H$ is nonempty, ample, and is closed under
coordinatewise intersection. Identify words with subsets of $E$, and put
\[
 \rho(v)=\{i\in v:v\setminus\{i\}\in H\}\quad(v\in H),\qquad
 \downarrow H=\{a:a\subseteq v\text{ for some }v\in H\}.
\]
For $a\in\downarrow H$ define its closure
$c(a)=\bigcap\{v\in H:a\subseteq v\}$. The intersection is nonempty
as a family, and belongs to $H$ by intersection closure.
The map $\rho$ is the established ordinary map on an intersection-closed
ample class. The construction below completes its native product basis
and supplies an explicit complementary map. It does not assume a basis
on the complement.

\begin{proposition}[The closure fibers are intervals]\label{prop:closure-intervals}
For every $v\in H$,
\[
 \{a\in\downarrow H:c(a)=v\}=[\rho(v),v].
\]
These intervals partition $\downarrow H$ and each contains exactly one
word of $H$, namely its top $v$.
\end{proposition}
\begin{proof}
If $c(a)=v$ and $i\in\rho(v)\setminus a$, the old word
$v\setminus\{i\}$ contains $a$, contradicting minimality of $c(a)$.
Thus $\rho(v)\subseteq a\subseteq v$.
Conversely suppose these containments hold and $c(a)\ne v$.
An ample class is isometric, so a Hamming geodesic from $c(a)$ to $v$
stays between those comparable words. Its penultimate word is
$v\setminus\{i\}\in H$ with $i\notin c(a)$. But then
$i\in\rho(v)\subseteq a\subseteq c(a)$, a contradiction.
Each $h\in H$ satisfies $c(h)=h$, proving the final assertion.
\end{proof}

For a convex geometry, the sets $\rho(v)$ are its extreme-point sets
and the interval description is its familiar unique-generator property.
The same proof works here without requiring either $\varnothing\in H$
or $E\in H$.
The published definitions of convex geometries and their bouquets are
recalled in Sections~2.4--2.5 of
\href{https://arxiv.org/abs/2405.12660}{Chalopin--Chepoi--Knauer,
\emph{Geometry of convex geometries}}. We use the elementary argument
above, not a new classification of that family.

\begin{theorem}[Complementary completion from closure intervals]
\label{thm:closure-completion}
For $K=W\setminus H$, define
\begin{equation}\label{eq:closure-complement-map}
 r_K(a)=
 \begin{cases}
 E\setminus a,&a\notin\downarrow H,\\
 (E\setminus c(a))\cup(a\setminus\rho(c(a))),&a\in\downarrow H\setminus H.
 \end{cases}
\end{equation}
This is an ordinary map made of actual whole cubes on the full original
complement. If the complement is nonempty, it has a corner peeling with the following order:
first process closure fibers in increasing $|v|$, deleting
$[\rho(v),v]\setminus\{v\}$ in decreasing word size; then delete
$W\setminus\downarrow H$ in increasing word size. Equal-size ties may
be broken arbitrarily. Its remaining incident directions at each deletion
are exactly \eqref{eq:closure-complement-map}.
Thus $H$ and its nonempty complement are corner-peelable, in particular $H$
belongs to $\mathrm{Ext}$. The construction is invariant under translating
the original coordinates and retains every original constant coordinate.
An empty complement has its empty map and order.
\end{theorem}
\begin{proof}
The whole $\rho(v)$-cube at $v$ lies in $H$: intersecting the lower
neighbors $v\setminus\{i\}$ gives each of its vertices. In any contained
cube its unique maximal word selects every free direction; each other
word has a zero in some free direction and omits it. Thus $\rho$ is
ordinary. Its normalized native vectors $e_{v,\rho(v)}$ form an
orthonormal basis of $V_H$ by the selected-basis criterion.

In the physical orthonormal coordinate basis from Section~\ref{sec:dual},
the support of $e_{v,\rho(v)}$ is exactly $[\rho(v),v]$.
Coordinates in $\rho(v)$ carry $(0,1)^t$, coordinates outside $v$ carry
$(1,0)^t$, and coordinates in $M=v\setminus\rho(v)$ carry $g_1$.
On this coordinate subspace use the full tensor basis obtained by choosing
$g_1$ or its orthogonal partner $e_{01}$ independently in $M$.
The all-$g_1$ vector is the existing primal vector. Each of the other
vectors has a unique root $a\in[\rho(v),v]\setminus\{v\}$ and is
\[
 e_{a,\,\rho(v)\cup(v\setminus a)}.
\]
At each coordinate index $a\notin\downarrow H$, use the standard tensor
vector $e_{a,a}$. Proposition~\ref{prop:closure-intervals} shows that
these constructions fill disjoint coordinate subspaces and use each
original root exactly once. They give an orthonormal basis of $V_H^\perp$
with roots precisely $K$.

Apply the orthogonal involution $U$ of Proposition~\ref{prop:dual}.
The resulting vectors have exactly the supports
\eqref{eq:closure-complement-map}. They span $V_K$, and
\eqref{eq:admissible-dual} verifies actuality of every whole selected cube.
Complement closure makes $K$ ample; the selected-basis criterion gives
every positioned native budget. This proves the ordinary assertion.

We verify the stronger peeling assertion directly. For $a$ of closure
$v$, its selected cube keeps $\rho(v)$ fixed at one, keeps
$v\setminus a$ fixed at zero, can remove bits of
$a\setminus\rho(v)$, and can add bits outside $v$.
With no outside bit added, all its other vertices are smaller words in
the same closure fiber and are deleted later. If an outside bit is added,
the closure, when defined, strictly contains $v$: for such a new word $w$,
intersect a containing old word with $v$ to see that it must contain
$c(w\cap v)=v$, since $w\cap v\in[\rho(v),v]$.
Otherwise the word lies outside $\downarrow H$, also processed later.
So the entire outgoing cube survives until $a$ is deleted.

A neighboring word in a direction omitted by $r_K(a)$ is already gone
or belongs to $H$. Removing $i\in\rho(v)$ gives a word with closure
contained in $v\setminus\{i\}\in H$, hence a strictly earlier fiber.
Adding $i\in v\setminus a$ gives a larger word in the same fiber,
or the excluded word $v$ itself. These are all omitted directions.
For $a\notin\downarrow H$, its remaining neighbors are exactly its
upward neighbors, and their whole cube survives by upward closure.
Thus every deleted word is a corner with the asserted incident cube.
The last surviving word may be retained instead of deleted.

The primal class peels by decreasing word size, as in the established
intersection-closed construction. None of the argument requires the zero
word. A translated intersection-closed class receives the construction
in that orientation; translating back transports both maps and orders
on the full original frame.
\end{proof}

The ordinary primal endpoint and Horn-renaming test are old. The relative
matrix conclusion here is the explicit complementary native basis and
its closure-fiber peeling order. Field-wide priority for this completion
formula is not certified. It includes the guarded-tip formula below:
there the nonsingleton closure fibers are exactly $[T_i,w_i]$.
By the complete nonnegative-product-basis classification in
\href{../../containment-cubic-positivity/research/local_nonnegative_projections.tex}
{\emph{Local nonnegativity of Boolean tensor projections}},
\texttt{thm:classification} and its parameter corollary, an ample class
whose projection is entrywise nonnegative in some real orthonormal product
basis has such an intersection-closed translate and therefore receives
this completion. This is a sufficient operator hypothesis; its existence
on arbitrary ample or cube-chordal classes is not assumed.

\section{A clamped cube attached to an intersection-closed source}
\label{sec:horn-gluing}
The preceding completion can be made compatible with an entire new cube.
The interface need not itself be a cube, or even a daisy. The construction
retains an arbitrary supplied primal map and supplies a complementary map;
it does not require a compatible old complementary map as input.

For a nonempty intersection-closed ample class $D$ on a coordinate set $J$,
write its full completion as
\begin{equation}\label{eq:full-closure-map}
 \Omega_D(x)=
 \begin{cases}
 \rho_D(c_D(x))\cup(c_D(x)\setminus x),&x\in\downarrow D,\\
 x,&x\notin\downarrow D.
 \end{cases}
\end{equation}
The proof of Theorem~\ref{thm:closure-completion} shows that
$(e_{x,\Omega_D(x)})_{x\in\{0,1\}^J}$ is an orthonormal basis of
the full physical space. At roots in $D$ it is the canonical primal
basis, since $\Omega_D(v)=\rho_D(v)$. These definitions also apply
when $J$ is empty: the nonempty class is then the singleton empty word.

\begin{theorem}[Compatible completion across a clamped cube]
\label{thm:horn-gluing}
Let $B\subseteq W$ be nonempty, ample and intersection-closed, let
$C=Q(z,S)$, and suppose the clamp into $C$ preserves $B$.
Then $H=B\cup C$ belongs to $\mathrm{Ext}$. One can retain every
column of any supplied ordinary map on $B$ and extend it to an ordinary
map on $H$, while constructing an ordinary map on $H^c$.
No daisy or chordality assumption is required.

Here is an explicit construction before an optional replacement of the
old primal map. Partition the original coordinates as
\[
 E=S\sqcup I\sqcup R,\qquad
 I=\{j\notin S:z_j=1\},\quad R=\{j\notin S:z_j=0\}.
\]
For $a\in F=\pi_SB$, set
\[
 L_a=\{y\subseteq I:(a,y,\varnothing_R)\in B\}.
\]
Each $L_a$, and $F$ itself, is nonempty, ample and intersection-closed.
With local supports embedded in their original coordinates, put
\begin{equation}\label{eq:horn-glued-basis}
 \Omega(x)=
 \begin{cases}
 \Omega_B(x),&x|_R\ne\varnothing,\\
 \Omega_F(a)\cup\Omega_{L_a}(y),
       &x=(a,y,\varnothing_R),\ a\in F,\\
 \Omega_F(a)\cup(I\setminus y),
       &x=(a,y,\varnothing_R),\ a\notin F.
 \end{cases}
\end{equation}
The vectors $(e_{x,\Omega(x)})_{x\in W}$ form a full native orthonormal
basis. The two ordinary
maps are $r_H(h)=\Omega(h)$ and $r_{H^c}(k)=E\setminus\Omega(k)$.
At every old root $b\in B$, $r_H(b)=\rho_B(b)$.
At every complementary root having a nonzero $R$-trace, the old canonical
complementary column is retained. Original constant coordinates remain
in $E$. A global coordinate flip transports the entire assertion, so it
also applies to any source having an intersection-closed orientation.
\end{theorem}
\begin{proof}
The proof splits the physical space at $R=0$, factors the old source
vectors there, and completes their remaining tensor factors.
For $b\in B$, the clamp and intersection closure give
\[
 b\setminus R=b\cap(b|_S,I,\varnothing_R)\in B.
\]
Thus setting all $R$-bits to zero also preserves $B$.
The old full basis \eqref{eq:full-closure-map} respects the physical
coordinate face $R=0$. To check this, if $x|_R=0$ and $c_B(x)$ exists,
zeroing $R$ in any old word above $x$ shows $c_B(x)|_R=0$.
Its entire interval $[\rho_B(c_B(x)),c_B(x)]$ lies in that face.
If $x|_R\ne0$ and $c_B(x)=v$, both $v\setminus R$ and $v$ belong
to $B$. A Hamming geodesic between them ends by adding some $j\in R$;
thus $j\in\rho_B(v)\subseteq x$. The full physical support interval
of this basis vector has that $j$ fixed at one, so misses $R=0$.
When closure does not exist the vector is the standard coordinate vector
at $x$, with the same conclusion. Consequently all old vectors at
roots with nonzero $R$-trace form a basis of the orthogonal complement
of the physical $R=0$ face. We keep them unchanged.

Projection closure of ampleness makes $F$ ample. Clamp invariance gives
the exact identity
\[
 F=\{a:(a,I,\varnothing_R)\in B\}.
\]
For each $a\in F$, the class $L_a$ contains $I$ and is a coordinate
section of $B$, so is ample. Projection and taking these sections
preserve intersection closure. Thus all the completions used in
\eqref{eq:horn-glued-basis} are defined, even if their classes omit zero.

The old canonical source vectors in the $R=0$ face factor as
\begin{equation}\label{eq:horn-source-factor}
 e_{(a,y,0),\rho_B(a,y,0)}
    = e_{a,\rho_F(a)}\otimes e_{y,\rho_{L_a}(y)}\otimes g_{0_R}.
\end{equation}
For an $S$-coordinate $j$, a lower neighbor $a\setminus\{j\}\in F$
gives the old word $(a\setminus\{j\},I,0)$; intersecting it with
$(a,y,0)$ gives that word's $j$-lower neighbor. Conversely, clamping
a $j$-lower neighbor of $(a,y,0)$ proves
$a\setminus\{j\}\in F$. Thus the $S$-part of $\rho_B$ is exactly
$\rho_F(a)$. Its $I$-part is $\rho_{L_a}(y)$ by the section definition,
and its $R$-part is empty, proving \eqref{eq:horn-source-factor}.

On the full physical $R=0$ face use the full $S$-basis for $F$.
Conditionally on its root $a\in F$, use the full $I$-basis for $L_a$.
For $a\notin F$, use instead the full $I$-tensor basis with local factors
$g_1$ at roots with bit one and $e_{01}$ at roots with bit zero.
Different $S$-roots are orthogonal independently of their $I$-factors;
at each fixed $S$-root the chosen $I$-vectors are a full orthonormal
basis. This proves that \eqref{eq:horn-glued-basis}, together with the
unchanged exterior, is a full orthonormal basis. Equation
\eqref{eq:horn-source-factor} proves literal retention of all old primal
vectors, including those outside $R=0$.

The new roots are exactly $(a,I,0)$ with $a\notin F$.
Their supports in \eqref{eq:horn-glued-basis} lie in $S$, so their
whole selected cubes lie in $C$. The old and new primal vectors are
therefore an orthonormal basis of $V_H$: they belong to that space
by \eqref{eq:cube-dictionary}, and there are $|H|$ of them.
All remaining vectors span $V_H^\perp$. Proposition~\ref{prop:dual}
turns them into actual complementary cube vectors with supports
$E\setminus\Omega(k)$. The old ample-union argument makes $H$ ample,
and complement closure makes $H^c$ ample. The selected-basis criterion
therefore pays every positioned cube equation on both shores.

Finally any supplied ordinary source basis spans the same space $V_B$.
Replace just the old canonical primal vectors by that basis, keeping
all new and complementary vectors. Orthonormality, actuality and all
conclusions persist. In particular cyclic supplied maps are allowed.
\end{proof}

\begin{corollary}[Simultaneous daisy pins with a marked source module]
\label{cor:horn-gluing-pins}
Under Theorem~\ref{thm:horn-gluing}, suppose in addition that $B\cap C$
is a daisy rooted at $d\in C$. There is an ordinary map on $B^c$ with
\[
 r_0(a)=E\setminus D(d,a)\qquad(a\in C\setminus B),
\]
retaining the complementary columns constructed by that theorem at every
root of $H^c$. The full basis before and after the attachment can be
related by a unimodular transformation over $\mathbb Z[q_i:i\in E]$
fixing the entire old vertex module on $B$ pointwise.
\end{corollary}
\begin{proof}
The new canonical daisy columns of Theorem~\ref{thm:relative}, normalized
by \eqref{eq:cube-dictionary}, form another orthonormal basis of
$V_H\cap V_B^\perp$. Replace only the new primal vectors by this bank.
Keep the entire old primal bank and every complementary vector.
Theorem~\ref{thm:excision} then supplies all the stated pins at once.

Before attachment use the full completion of $B$; after attachment use
the full basis just constructed, with the complementary supports
dualized back to kernel supports. Both are native bases of the full
cube and therefore give polynomial actual-cube matrices $P_0,P_1$
of determinant one, by the first-paper integral selected-basis theorem.
They have identical columns at roots of $B$. Those columns form a
unimodular basis of the vertex module on $B$, so
$P_1P_0^{-1}$ fixes that module pointwise. If a different supplied primal
map is retained, replace the old bank in both full bases first.
\end{proof}

The source need not remain intersection-closed after this operation.
For example, the old atlas class $V$ in the next subsection is obtained by
one such attachment to a Horn-oriented remainder, but neither shore of
$V$ has a Horn orientation. Thus this theorem does not provide an
iteration for arbitrary cube-chordal classes. It extends the earlier
down-set attachment theorem to arbitrary intersection-closed ample sources,
and supplies the full jointly prescribed bank rather than independently
admissible columns. The existence domain of every particular example
may already follow from other constructions; no field-wide priority
claim is made for this gluing statement.

The exact verifier \texttt{verify\_horn\_cube\_gluing.py} performs 4993
positioned-equation checks in six fixed controls. They include the old atlas
attachment, a five-pin noncube interface on six original coordinates,
a non-daisy interface, a source omitting zero, a cyclic supplied source
and an attachment adding no words. These are checks of the explicit
formulas, not substitutes for the proof or a search for a separation.

\subsection{Both shores can lack a nonnegative product basis}
\label{sec:two-shore-horn-obstruction}
Use the existing atlas class, in low-bit-first binary masks,
\[
 V=\{0,1,2,3,6,7,8,9,16,17,18,19,21,23,24,26\}\subseteq\{0,1\}^5.
\]
Its ampleness and chordal maximal-cube graph are proved in the atlas
\href{../../minor-atlas/sections/cube_chordal_ample.tex}
{example on projection and complementation}.
For clarity its complete maximal cubes are
\[
 (0,19),(2,5),(17,6),(0,9),(16,10),(3,20),(0,24).
\]
The following prime missing patterns suffice to rule out Horn renaming
on each shore; columns are support mask and missing-pattern mask:
\[
\begin{array}{c|cc}
 S& V& W\setminus V\\\hline
26&10&2\\
22&4&0\\
21&20&16\\
25&25&17
\end{array}
\]
In the established Horn test (\texttt{lem:horn-test} of that same original),
the first two rows forbid respectively
$(r_1,r_4)=(1,0)$ and $(0,0)$, forcing $r_4=1$.
The last two forbid $(r_0,r_4)=(0,1)$ and $(1,1)$, forcing $r_4=0$.
This contradiction applies to both pattern columns. Thus neither $V$
nor its full complement has an intersection-closed coordinate translate.
CP's classification therefore excludes entrywise nonnegativity of either
projection in any real orthonormal product basis. This applies the received
CP theorem to a smaller CAmp host; it is not a new spectral classification
or a negative ordinary-map example.

\section{Reselecting a complementary basis across a new interface}
\label{sec:reselecting}
The same class $V$ tests why the old complementary map must be allowed
to change. It has four simplicial maximal cubes. For each, remove its
private word $a$ and write $B=V\setminus\{a\}$. The interface is a
three-word path with its unique daisy root $z$. Every such $B$ has
intersection-closed coordinate translates. Nevertheless none of the
canonical closure completions, over any of its possible centers, satisfies
the one prescribed support needed to add that cube:
\[
\begin{array}{c|r|r|c|r|r}
 C &a&z&\text{all centers of }B&r_{B^c}(a)&E\setminus D(z,a)\\\hline
 (0,9)&9&0&17,19&6&22\\
 (2,5)&6&3&16,18&10&26\\
 (16,10)&26&16&2,3&5&21\\
 (17,6)&21&19&0,1&9&25
\end{array}
\]
The canonical values in the fifth column are the same at both listed
centers. The complete finite check tests all original five-bit centers,
not a chosen center. Thus even a Horn old host does not make its canonical
closure completion compatible with an arbitrary simplicial attachment.
The assertion concerns this canonical family, not all maps on $B^c$.

\subsection{A standard phase move retains the old primal space}
Let $p$ be any ordinary map on an ample $B$, and let $r$ be any ordinary
map on its full complement. Their ambient outmap is
\[
 O(v)=p(v)\quad(v\in B),\qquad O(v)=E\setminus r(v)\quad(v\notin B).
\]
In the native product dictionary its vectors form an orthonormal basis
of the full tensor space: the second bank is the $U$-image of the
complementary bank and spans $V_B^\perp$. Hence $O$ is a full-cube USO.
This is the established complementary paste, used with real supplied maps.

Fix a coordinate $i$. Form a graph on the original words, joining $u,v$
when $u_i\ne v_i$ and
\begin{equation}\label{eq:phase-graph}
 (D(u,v)\setminus\{i\})\cap(O(u)\mathbin\triangle O(v))=\varnothing.
\end{equation}
Its components are the endpoint sets of the standard $i$-phases.
Toggling the $i$-bit of $O$ at every word of a component preserves
orthogonality: a pair whose orthogonality relies only on $i$ is connected
by \eqref{eq:phase-graph} and has both bits toggled; every other pair
retains an orthogonal outside factor. Every component contains both
endpoints of each of its $i$-edges. This is the classical phase-flip
operation, not a new exchange primitive; see Definition~2.4 and
Observation~2.1 of Borzechowski--Weber,
\href{https://doi.org/10.1007/s00373-025-02929-2}
{\emph{On Flipping Edge Sets in Unique Sink Orientations}}.

If the component avoids $B$, all old primal vectors remain literal and
still span $V_B$. Its other vectors therefore still span $V_B^\perp$.
Applying $U$ back gives an ordinary map on $B^c$, with only those
$i$-support bits changed and every selected cube still actual. This verifies
the marked-space contract of the old phase move. It does not assert that
a component avoiding $B$ is always available for a desired prescription.

There is also an integral polynomial formulation. Let $P_O,P_{O'}$ be
the full-cube matrices of unnormalized selected cube columns before and
after the move. Both are unimodular over $\mathbb Z[\mathbf q]$ by the
selected-basis theorem. Thus $P_{O'}P_O^{-1}$ is a polynomial unimodular
operator carrying each old selected column to its new actual cube column.
It fixes the entire vertex module on $B$ pointwise: the unchanged primal
columns are already a polynomial unimodular basis of that module.
This preserves the marked source module and the selected rooted-cube data;
invariance of every other positioned-cube subspace is not asserted.

For the four rows of the preceding table, choose respectively centers
$17,16,2,0$. A single phase in coordinate $4$ avoids $B$ and repairs the
prescription. Its endpoint sets are respectively
\[
 \{9,25\},\quad\{4,6,20,22\},\quad
 \{10,26\},\quad\{4,5,20,21\}.
\]
Every old primal column and every complementary column outside that phase
is retained. The repaired complement map restricts unchanged to $V^c$
by Theorem~\ref{thm:excision}. This is an explicit source-preserving
repair on the fixed example, not a universal phase-supply theorem.

\subsection{Choosing between clique-bag roots}
Another restricted alphabet comes from the
\href{../../maxmin-forbidden/sections/chordal_downset_tree.tex}
{positioned clique-tree theorem} and its proposition on opposite bag roots.
Fix a root $z_t$ in every clique bag. For $a\in K=H^c$ let
\[
 A(a)=\{t:Q(a,E\setminus D(a,z_t))\subseteq K\}.
\]
Every $A(a)$ is nonempty. Indeed choose a minimally missing pattern
matched by $a$; the opposite-root proposition supplies a bag root opposite
its entire trace. The corresponding cube keeps that forbidden trace fixed.
This is reuse of the complete positioned-support theorem, not a consequence
of abstract graph chordality alone.

Selecting $t(a)\in A(a)$ gives an ordinary map from these particular
columns exactly when
\begin{equation}\label{eq:chart-compatibility}
 D(a,b)\setminus D(z_{t(a)},z_{t(b)})\ne\varnothing\qquad(a\ne b).
\end{equation}
To see this, two native factors with different root bits are orthogonal
exactly when their support bits differ. For the displayed away-from-root
supports, this happens precisely where the two bag roots agree. All other
local overlaps are nonzero at interior parameters. Thus
\eqref{eq:chart-compatibility} is exactly pairwise orthogonality; the
selected-basis theorem gives the assertion on the ample class $K$.
Nonempty domains alone do not give a simultaneous selection. This
restricted alphabet is sufficient; Theorem~\ref{thm:camp-ext} now
constructs a simultaneous selection by tree completion. Necessity of
this particular alphabet is not asserted.

For the fixed $V$, the four bag roots are $0,3,16,19$.
One compatible selection uses the following root groups:
\[
\begin{array}{c|l}
 z&\text{roots }a\in V^c\text{ using }z\\\hline
0&25\\
3&22\\
16&10,11,27\\
19&4,5,12,13,14,15,20,28,29,30,31
\end{array}
\]
Each selected support is $E\setminus D(a,z)$.
The actual cubes, condition \eqref{eq:chart-compatibility}, and all native
rows are checked directly in the owned certificate. This is a positive
example outside the two-shore Horn sufficient condition, not a new claim
about its prior membership or a proof of the universal selection.

A fixed priority of bag roots is insufficient even on a median example.
Use the existing arbitrary-root-pattern construction on a four-node path,
with distinguished bits $0,1,0,1$. Its twelve words on eight coordinates are
\[
 \{0,1,3,7,8,9,11,15,16,41,67,143\}.
\]
The companion's complete arbitrary-root-pattern theorem proves it is median
and CAmp, with forced bag roots $0,9,3,15$ along its unique clique-tree path.
If each complementary word takes its first eligible root in this order,
word $64$ takes root $9$ and word $72$ takes root $0$.
Their actual supports are $182$ and $183$. Neither includes direction
mask $8$, so the actual edge $\{64,72\}$ has budget zero.
The whole restricted alphabet still admits a map on this fixed example;
the diagnostic only rejects the priority rule. The general compatible
selection in Theorem~\ref{thm:camp-ext} uses the stages of tree completion,
which are determined by joint geometry rather than individual eligibility.

\section{Normalization along a monotone coordinate block}
\label{sec:monotone-normalization}
Here is a situation in which all the prescribed supports can be supplied
at once, retaining the induced maps on the other parallel faces. Its
ordinary construction is the established nested-coordinate lift, iterated
over a block. The full original proof at
\href{../../containment-ext-rep/sections/nested_coordinate.tex}{the nested-coordinate module},
\texttt{thm:er-nested-selector}, supplies the one-coordinate predecessor.
The additional description here is its action on the precise simultaneous
prescriptions of Theorem~\ref{thm:excision}.

Reflect the $Y$-bits so that $z|_Y=0$. Say that $B$ is downward closed
in the block $Y$ if $(u,y)\in B$ implies $(u,y')\in B$ whenever
$y'\subseteq y$. For $K_0=B^c$, let
$(K_0)_y=\{u:(u,y)\in K_0\}$. These sections increase with $y$.

\begin{proposition}[Idempotent normalization with prescribed supports]
\label{prop:monotone-retraction}
Assume $Y\ne\varnothing$, the hypotheses of Theorem~\ref{thm:excision},
and downward closure of $B$ in $Y$. From any ordinary map $r_0$ of
$K_0$, construct on $K=K_0\setminus C$ the map
\begin{equation}\label{eq:monotone-new}
 r(u,y)=\bigl(r_0(u,y)\cap S\bigr)
               \cup\{j\in Y:y_j=0\},\qquad y\ne0.
\end{equation}
Append the supports $E\setminus T_a$ at $a\in A$.
This gives an ordinary map $N(r_0)$ on $K_0$ satisfying every prescription
in \eqref{eq:pins}. It retains the induced ordinary map on every parallel
$S$-face with $y\ne0$, and $N^2=N$. Cycles in those face maps are retained.
\end{proposition}
\begin{proof}
Conditioning $r_0$ gives an ordinary map $r_0(u,y)\cap S$ on each
section $(K_0)_y$. Their whole cubes remain actual when $Y$-coordinates
are increased, by upward closure of $K_0$ in $Y$. The additional
directions in \eqref{eq:monotone-new} only increase zero bits and cannot
reach $y=0$ from a nonzero $y$. Thus every selected cube lies in $K$.

Write a positioned cube $Q\subseteq K$ as $Q_S\times Q_Y$.
The coordinate orientation on $Y$ in \eqref{eq:monotone-new} has its
unique budget owner at the minimum word $y_*$ of $Q_Y$. This word
is nonzero, since otherwise $Q$ would meet the deleted face.
The face $Q_S$ lies in $(K_0)_{y_*}$, where its ordinary conditioned
map supplies exactly one owner. This checks every cube row of $K$.
Theorem~\ref{thm:excision} supplies the appended map on $K_0$.
Formula \eqref{eq:monotone-new} leaves each retained $S$-restriction
unchanged and replaces only its $Y$-directions; a second application
therefore has no further effect. The appended supports are fixed too.
\end{proof}

\section{Simultaneous prescribed supports above a down-set}
\label{sec:guarded-tips}
The next construction handles several attachments together even when
the final host has no nested coordinate. It is a simultaneous application
of the established full-collar filling construction, whose complete proof
is \texttt{prop:damp-collar-independent-maps} in the
\href{../../dismantlable-ample/sections/classification/local_extenture_shortening.tex}
{local shortening module}. The distinction here is the explicit retained
columns and all the canonical prescriptions, rather than a new generic
filling operation or an unrestricted existence theorem.

Identify words with subsets of $E$.

\begin{theorem}[One cube attached to a down-set]\label{thm:downset-attachment}
Let $D\subseteq2^E$ be a nonempty down-set, let $C=Q(z,S)$, and
assume its clamp preserves $D$. Then $D\cup C\in\mathrm{CAmp}\cap\mathrm{Ext}$.
More precisely choose $z|_S=0$ without changing $C$, and partition
$Y=E\setminus S$ into $I=\{j:z_j=1\}$ and $R=\{j:z_j=0\}$.
Put $F_S=\{u:(u,I,0_R)\in D\}$ and $A_S=2^S\setminus F_S$.
On $D^c$ the following ordinary map satisfies all the canonical
prescriptions at the new roots $a\in C\setminus D$:
\begin{equation}\label{eq:downset-attachment-map}
 r_0(v)=
 \begin{cases}
 (S\setminus v)\cup(v\cap I)\cup R,&v\cap R=\varnothing,\\
 E\setminus v,&v\cap R\ne\varnothing.
 \end{cases}
\end{equation}
Its unchanged restriction to $(D\cup C)^c$ is ordinary. In particular
it retains every canonical upward column at roots with a nonzero $R$-bit.
The interface $F_S$ may be any down-set, and the construction uses no
chordality or bound on dimension.
\end{theorem}
\begin{proof}
The clamp of the zero word belongs to $D$, so $F_S$ is nonempty and
downward closed. Clamp invariance and downward closure give the exact
identity
\[
 D[R=0]=F_S\times2^I.
\]
Indeed membership on the right follows by lowering the $I$-bits in a
clamped word, and membership on the left implies its $S$-trace belongs
to $F_S$ by clamping. Moreover any $S$-trace in $A_S$ is absent from
$D$ at every $Y$-assignment, again by the clamp.

Let $K_* = D^c[R\ne0]$. It is an up-set, so its canonical upward map
is ordinary. The full $R$-collar over $A_S\times2^I$ lies in $K_*$.
When $R\ne\varnothing$, apply the established collar filling to this
seed and either of the two core classes
\[
 L_0=A_S\times2^I,\qquad
 L=A_S\times(2^I\setminus\{I\}).
\]
Use the product of the upward map on the up-set $A_S$ and the downward
map on the indicated $I$-class. Both are ordinary: in the punctured
$I$-cube a proper subset remains proper when bits are lowered, and every
contained face has its unique maximal owner. The filling adds all $R$
directions and gives exactly \eqref{eq:downset-attachment-map}, first
on $D^c$ and then on $(D\cup C)^c$. The columns on $K_*$ agree.
If $R$ is empty, there is no seed and these same two product maps give
the conclusion directly. If $I$ is empty, $L$ is empty and the new
complement is just $K_*$, with its upward map. Empty cores use the
empty-map convention; if $A_S$ is empty there is no change at all.

For a new root $a=(u,I,0_R)$, the support in
\eqref{eq:downset-attachment-map} is
$(S\setminus u)\cup I\cup R=E\setminus D(z,a)$.
The union is ample by the atlas clamp argument. The ordinary downward
map on $D$ extends to it by Theorem~\ref{thm:relative} and the selected
cube-basis criterion. Thus both required ordinary maps exist.
For cube chordality, if $I$ is empty the union is still a down-set.
Otherwise $C$ contains no old maximal cube, since every old maximal cube
contains zero and $C$ does not. The old maximal-cube graph is complete,
and the new cube is a simplicial vertex in it. The atlas generator gives
cube chordality. The case $C\subseteq D$ makes no change.
\end{proof}

The operation in this proof is the previously established full-collar
filling, with an explicit identification of its seed, two core classes
and prescribed old-complement columns. The common interface $F_S$ need
not be a cube, so it is not an application of a common-cube assembly
theorem without this geometric verification. No external priority claim
for this derived attachment criterion is made.
Its ordinary existence also follows from the established nested-coordinate
recursion: first split the $R$-coordinates, directed downward, whose one
sections are down-sets. In the remaining $R=0$ face, split the
$I$-coordinates, directed upward; a zero section is a product with $F_S$
and the final all-one section is a full $S$-cube. Thus this construction
does not enlarge that old existence domain. Its added contract is the
explicit simultaneous prescriptions and retained exterior columns.

Now let $D\subseteq2^E$ be a nonempty down-set, and choose finitely many
distinct words $w_i$ and supports
$T_i\subseteq w_i$ such that
\begin{equation}\label{eq:guarded-conditions}
 T_i\notin D,\qquad
 \{v\subseteq w_i:T_i\nsubseteq v\}\subseteq D,\qquad
 T_i\nsubseteq w_j\quad(i\ne j).
\end{equation}
Put $I_i=w_i\setminus T_i$, $O_i=[T_i,w_i]$ and
$H=D\cup\{w_i:i\}$. The intervals $O_i$ are pairwise disjoint,
and are wholly outside $D$.

\begin{theorem}[Simultaneous local normalization]\label{thm:guarded-pins}
The class $H$ is ample and has the ordinary map
\[
 r_H(v)=v\quad(v\in D),\qquad r_H(w_i)=T_i.
\]
On the old complement $D^c$ the formula
\begin{equation}\label{eq:guarded-map}
 r_0(v)=
 \begin{cases}
 (v\cap I_i)\cup(E\setminus w_i),&v\in O_i,\\
 E\setminus v,&v\notin\bigcup_i O_i
 \end{cases}
\end{equation}
is an ordinary map satisfying simultaneously
$r_0(w_i)=E\setminus T_i$. It retains the canonical upward map at
every root outside the displayed intervals. Its unchanged restriction
to $H^c$ is ordinary. In particular $H\in\mathrm{Ext}$.
\end{theorem}
\begin{proof}
For ampleness, add each cube $C_i=Q(w_i,T_i)$ in any order.
It adds just $w_i$: every other word of $C_i$ is an old subset of
$w_i$ missing part of $T_i$, hence belongs to $D$.
Every old word has a proper $T_i$-trace, since $T_i\notin D$ and
$T_i\nsubseteq w_j$ for $j\ne i$. Its clamp into $C_i$ therefore lies
in $D$ by \eqref{eq:guarded-conditions}. The atlas ample-union argument
applies at every step. Reversing these corner additions gives the
displayed ordinary primal map, by the previously reviewed corner extension.

For the opposite construction, first consider a root outside all $O_i$.
Its upward cube stays outside $D$. It cannot contain a deleted word $w_i$:
that would mean $v\subseteq w_i$, and exclusion from $D$ would then
force $T_i\subseteq v$, placing $v$ in $O_i$.
For $v\in O_i$, the cube in \eqref{eq:guarded-map} keeps $T_i$ fixed
at one, decreases bits of $I_i$, and increases bits outside $w_i$.
It avoids $D$ because $T_i\notin D$, and avoids every $w_j$ with
$j\ne i$ because $T_i\nsubseteq w_j$. If $v\ne w_i$, some bit of
$I_i$ is fixed at zero, so it also avoids $w_i$. This proves whole-column
actuality both on $D^c$ and, after deletion, on $H^c$.

We check every cube row on either host. If a cube meets no $O_i$,
all its roots have the upward map, whose unique owner is its minimum word.
Otherwise take $v\in Q\cap O_i$. No direction of $T_i$ is free in $Q$:
flipping it at $v$ would give a subset of $w_i$ lacking $T_i$, hence a
word of $D$. Thus $T_i$ is fixed at one throughout $Q$, which also
prevents $Q$ from meeting any other $O_j$. Every fixed coordinate outside
$w_i$ is zero. A root of $Q$ with a positive bit outside $w_i$ is therefore
an unchanged root missing that free direction from its support, and cannot
own $Q$. At outside bits zero, the remaining row is a face on $I_i$,
with $T_i$ fixed at one. The map on this face selects its positive bits,
so its unique owner is the maximum word of that face. On $H^c$ the face
excludes the full word $I_i$; this orientation remains an actual ordinary
map of that punctured cube. All outside directions belong to
\eqref{eq:guarded-map}, so the same unique root owns the full row $Q$.
Finally $r_0(w_i)=I_i\cup(E\setminus w_i)=E\setminus T_i$, as required.
\end{proof}

This entire existence domain already lies in the physical square-free
shore domain stated in ER109, whose supplier is credited there to FA035.
Indeed the $w_i$ form an inclusion antichain by
\eqref{eq:guarded-conditions}. At any coordinate $c$, the exclusive
one-shore consists exactly of $w_i\setminus\{c\}$ for those $c\in I_i$.
It is again an antichain and contains no edge, hence no square.
The difference supplied here is the explicit joint canonical pins and
retention of all other upward columns, not first membership in that domain.

\subsection{A chordal sun requires joint normalization}
\label{sec:sun-normalization}
We use the existing chordal sun family, not a new peripheral obstruction.
Its complete original statement and proof are in
\href{../../cover-recursive/research/rper_chordal_sun_obstructions.tex}
{the peripheral companion}, \texttt{thm:rper-chordal-sun-obstructions}.
For $m\ge4$, indices are cyclic and
\[
 E=\{(i,i-1),(i,i+1):i\in\mathbb Z_m\},\quad
 A_i=\{(i-1,i),(i+1,i)\},\quad B_i=\{(i,i-1),(i,i+1)\},
\]
\[
 D=\bigcup_i2^{A_i\cup A_{i+1}},\qquad
 w_i=A_i\cup B_i,\qquad H=D\cup\{w_i:i\in\mathbb Z_m\}.
\]
The original proof shows $H\in\mathrm{CAmp}$, $|H|=13m+1$, VC dimension
four and no peripheral coordinate. Its maximal cubes are the core
four-cubes and the pairwise disjoint leaf squares $Q(w_i,B_i)$.
The conditions \eqref{eq:guarded-conditions} hold with $T_i=B_i$ and
$I_i=A_i$: missing either outgoing bit gives an old core word, while no
other tip contains both outgoing bits. Thus Theorem~\ref{thm:guarded-pins}
supplies all canonical supports $E\setminus B_i$ in one ordinary old
complement map, and gives an explicit ordinary map on $H^c$.

There is no new pointwise existence domain claimed for these suns.
For a coordinate $c=(f,e)$, the physical exclusive shore
$H[c=1]\setminus H[c=0]$ is the singleton $w_e\setminus\{c\}$.
All old down-set words have their lower neighbor; among the two tips
using $c$, only $w_e$ lacks it. Thus the family already satisfies the
square-free physical shore hypothesis of
\href{../../containment-ext-rep/research/ER109_SQUARE_FREE_PRIVATE_COMPLEMENTS.md}
{ER109}, Theorem~4, whose existence supplier is credited there to FA035.
We inspected the complete ER109 argument, retaining its AUTHOR status and
its dependency scope. Our direct formula is a separate proof on this
family with the additional explicit retained-column and simultaneous-pin
contract. We do not receive its general supplier by checking these suns.

This example also prevents a coordinate-wise commutation shortcut.
Fix $i$ and take the actual attachment
$B=H\setminus\{w_i\}$, $C=Q(w_i,B_i)$ with root $z=A_i$.
The full clamp into $C$ preserves $B$. Nevertheless, for every coordinate
$c\in E\setminus B_i$, clamping that coordinate alone toward $z_c$
fails to preserve $B$. If $c\in A_i$, its target bit is one; add it
to a core four-set not using target $i$. The result has five bits and
is absent from $B$. Otherwise $c=(f,e)$ with $e\ne i$ and target bit
zero. The retained word $w_e$ loses an incoming bit under this clamp;
the result still contains $B_e$, so it is absent from $D$ and, having
three bits, from every tip. These cases exhaust $E\setminus B_i$.
Proposition~\ref{prop:commuting-clamp} therefore says that $\Pi_B$
commutes with the whole $\Pi_C$, but with none of its one-coordinate
factors. Independent coordinate normalization is not justified by the
whole-cube commutation, even on a genuine atlas attachment.

The general problem now asks for joint normalization across successive
daisy interfaces with changing roots. The monotone-block and guarded-tip
constructions pay this compatibility only at their stated hypotheses.

\section{What summing the inverse diagonal forgets}
\label{sec:rooted-trace}
For an ample class $D$, put $x_i=q_i^2/(1-q_i^2)$ and
$y_i=1+2x_i$. The first-paper inverse enumerator gives
\begin{equation}\label{eq:rooted-diagonal}
 (G_D^{-1})_{vv}=L_v(\mathbf x),\qquad
 L_v(\mathbf x)=\sum_{T:Q(v,T)\subseteq D}x_T.
\end{equation}
Summing over roots gives
$\operatorname{tr}G_D^{-1}=\sum_{T\in\Sh(D)}y_T$.
The complete labelled trace therefore determines exactly the shattered
family; the equal-parameter trace determines only its numbers in each size.
These facts and their median, maximum and coordinate-product applications
are proved in the
\href{../main.pdf}{completed first paper}, in the sections on inverse
enumeration and trace subclasses. They are reused here.

The individual diagonal polynomials retain the positions lost by summation:
the coefficient of $x_T$ in $L_v$ is exactly the indicator that $Q(v,T)$
is an actual cube. Thus the family $(L_v)_{v\in D}$ recovers all actual
positioned cubes. This is an immediate consequence of the existing inverse
formula, not an independent existence theorem for cube bases.

\begin{example}[The trace does not detect cube chordality]
Let $P=\{000,100,110,111\}$ be a three-edge path and
$Z=\{000,100,010,001\}$ a three-leaf star. Both are ample and have
shattered family consisting of the empty support and the three singletons.
On two disjoint triples of labelled coordinates, the ample classes
$P\times P$ and $Z\times Z$ therefore have the identical trace
\[
 (1+y_1+y_2+y_3)(1+y_4+y_5+y_6).
\]
Every maximal cube of $Z\times Z$ is a product of two star edges and
contains $(000,000)$. Its maximal-cube intersection graph is complete,
so it is cube-chordal. Number the consecutive edges of $P$ by
$E_0,E_1,E_2$. In $P\times P$ the four maximal squares
\[
 E_0\times E_1,\quad E_1\times E_0,\quad
 E_2\times E_1,\quad E_1\times E_2
\]
form an induced four-cycle in that order: consecutive squares meet at
one word, while each opposite pair is disjoint. Thus $P\times P$ is
not cube-chordal. This is a trace application of the atlas product
criterion, \texttt{prop:cube-chordal-product}; it is not a new
classification of cube-chordal products and says nothing negative about
ordinary complementary existence on either host.
\end{example}

For example, let $N_D(v)=\{i:v^{(i)}\in D\}$,
where $v^{(i)}$ is $v$ with coordinate $i$ flipped. A vertex is a
\emph{corner} if it belongs to a unique maximal positioned cube.
Every rooted cube at $v$ uses directions in $N_D(v)$, so
\begin{equation}\label{eq:corner-diagonal}
 (G_D^{-1})_{vv}\le\prod_{i\in N_D(v)}(1+x_i)
                   =\frac{1}{t_{N_D(v)}}.
\end{equation}
Equality at one interior parameter point holds exactly when $v$ is a corner.
Indeed the difference is the positive sum of $x_T$ over the missing
subsets of $N_D(v)$. Equality means the whole neighbor cube is actual,
which contains every rooted cube at $v$ and is the unique maximal one.
Conversely every neighbor edge extends to a maximal cube containing $v$;
uniqueness forces all neighbor directions into that cube. This is the
known corner criterion, expressed through a marked inverse entry.

One can apply the test successively without inverting each smaller Gram
matrix from scratch. Write $K=G_D^{-1}$ and $J=D\setminus\{v\}$. Then
\[
 G_J^{-1}=K_{JJ}-\frac{K_{Jv}K_{vJ}}{K_{vv}},
\]
by the block inverse identity. Corner deletion preserves ampleness, so the
next marked diagonal test applies. A sequence of successful deletions down
to one vertex certifies a corner peeling. We assert no completeness of an
arbitrary greedy choice, and no implication from failure of this procedure
to absence of an ordinary map. The full trace alone cannot determine
corner peelability, as the first paper's same-shadow examples show.
\end{document}
